\documentclass[12pt]{ctexrep}

\usepackage[a4paper,margin=2.5cm]{geometry}
\usepackage{setspace}
\usepackage{graphicx}
\usepackage{booktabs}
\usepackage{longtable}
\usepackage{array}
\usepackage{amsmath}
\usepackage{hyperref}
\usepackage{float}

\hypersetup{
    colorlinks=true,
    linkcolor=black,
    urlcolor=blue,
    citecolor=black
}

\setstretch{1.25}
\setlength{\parindent}{2em}
\setlength{\parskip}{0.3em}

\title{基于多模态信息融合的驾驶员分心与疲劳识别方法研究}
\author{黄景泓}
\date{2026年4月}

\begin{document}

\maketitle
\tableofcontents
\clearpage

\chapter{绪论}

\section{研究背景}
随着智能交通和辅助驾驶技术的快速发展，驾驶员状态识别逐渐成为车内感知研究中的重要方向。在实际驾驶过程中，驾驶员可能因接打电话、饮水、操作中控、与乘客交流等行为产生分心，也可能因长时间驾驶而出现打哈欠、闭眼、困倦等疲劳状态。这些异常状态均会显著提高交通事故发生风险，因此构建面向真实车内场景的驾驶员状态识别方法具有重要的工程应用价值。

目前公开数据集上的驾驶状态识别模型通常可以取得较高精度，但当模型迁移到真实车内场景时，往往会受到拍摄视角、光照条件、车内布局、遮挡程度以及标签定义差异等因素影响，导致性能明显下降。尤其是疲劳识别任务，其对面部细粒度特征高度敏感，跨数据集泛化难度更高。因此，如何在公开数据集训练能力的基础上，通过目标域适配以及多模态融合提升模型在真实车内环境中的表现，是本文重点研究的问题。

\section{研究目标与主要问题}
本文的总体目标是构建一个能够在真实车内场景下稳定运行的驾驶员状态识别系统，实现对驾驶员分心状态和疲劳状态的联合识别。围绕这一目标，本文主要解决以下几个问题：

\begin{enumerate}
    \item 如何利用公开数据集建立有效的单模态分心识别模型，并将其迁移到真实车内场景中。
    \item 如何以面部信息为主完成疲劳识别，并解决疲劳模型在不同数据集之间存在显著域差的问题。
    \item 如何设计 body 与 face 双模态多任务模型，利用不同模态之间的互补信息同时提升分心识别与疲劳识别性能。
    \item 如何将训练完成的模型进一步实现为可运行的本地推理与可视化系统，为论文答辩演示与后续应用提供支持。
\end{enumerate}

\section{本文的主要工作}
围绕上述问题，本文目前已完成的主要工作如下：

\begin{enumerate}
    \item 构建了以 State Farm 为源域、DMD S1 为目标域的分心识别流程，完成了单模态 body 分类模型训练、目标域混合微调以及误判分析，并确定了当前较优的分心单模态模型。
    \item 构建了以 NTHU 为源域、DMD S5 为目标域的疲劳识别流程，完成了基于 face 模态的疲劳二分类模型训练、源域到目标域的迁移适配以及误差分析，并确定了当前较优的疲劳单模态模型。
    \item 设计并实现了 body+face 双模态多任务模型，通过共享双模态特征表示联合完成分心识别与疲劳识别，并在统一测试集上完成了单模态与多模态的公平对比。
    \item 实现了双视频对齐脚本、多模态视频推理脚本以及本地界面原型，为模型可视化展示与系统集成奠定了基础。
\end{enumerate}

\section{论文结构安排}
本文共分为八章。第一章介绍课题背景、研究目标和主要工作；第二章介绍相关理论与关键技术；第三章给出系统总体方案设计；第四章和第五章分别介绍单模态分心识别与单模态疲劳识别方法；第六章介绍双模态多任务模型设计与实验结果；第七章介绍系统实现与界面设计；第八章总结全文并展望后续工作。

\chapter{相关理论与关键技术}

\section{卷积神经网络与图像分类}
卷积神经网络是视觉分类任务中最常用的基础模型之一。其通过卷积层、池化层和非线性激活函数逐步提取图像特征，能够从底层边缘纹理逐渐学习到高层语义表示。对于驾驶员分心识别与疲劳识别而言，无论输入为 body 图像还是 face 图像，本质上都属于基于图像特征的分类问题，因此卷积神经网络仍然是建立基础识别模型的重要手段。

\section{迁移学习与微调}
迁移学习是指利用模型在大规模源域数据上学到的通用视觉特征，并通过微调迁移到目标任务中。对于本文的任务而言，公开数据集与真实车内数据在分布上存在明显差异，因此单纯依赖源域训练通常难以获得理想的目标域性能。本文在分心任务中采用了 State Farm 到 DMD S1 的迁移微调策略，在疲劳任务中采用了 NTHU 到 DMD S5 的迁移微调策略，以提升真实车内场景下的识别效果。

\section{多模态融合思想}
多模态学习的核心思想是联合利用不同来源的信息进行建模。对于本文研究对象，body 模态更擅长表达驾驶姿态、手部动作和整体行为，face 模态更擅长表达眼睛状态、嘴部张开和疲劳表情，两种模态具有天然互补性。因此，构建 body+face 双模态模型有望弥补单模态信息不足的问题。

\section{多任务学习与 masked loss}
多任务学习是指在统一模型中同时优化多个相关任务。本文的多模态模型同时输出分心识别结果和疲劳识别结果，但并非每一条样本都同时具备完整的分心标签和疲劳标签，因此在训练时采用 masked loss 策略：当样本缺少某一任务标签时，对应任务损失不参与该样本的反向传播。其基本形式可以表示为：

\begin{equation}
L = \lambda_d L_{distraction} + \lambda_f L_{fatigue}
\end{equation}

其中 $L_{distraction}$ 和 $L_{fatigue}$ 分别表示分心任务与疲劳任务的分类损失，$\lambda_d$ 和 $\lambda_f$ 为任务权重系数。该设计使得不同来源、不同标注完整度的数据能够在统一框架中联合训练。

\chapter{系统总体方案设计}

\section{总体流程}
本文系统整体采用“单模态建模验证 + 目标域适配 + 双模态融合 + 本地化实现”的技术路线。首先分别完成分心识别与疲劳识别的单模态建模，明确各自的基础性能与目标域问题；随后构建 body+face 双模态多任务模型，验证多模态互补性；最后在训练完成的模型基础上实现视频推理与本地界面原型。

\section{系统模块划分}
为便于后续实验组织与系统实现，本文将整体工作划分为以下模块：

\begin{enumerate}
    \item 数据处理模块：完成多数据集的清洗、标签映射、训练集划分与配对构建。
    \item 单模态训练模块：分别完成分心单模态模型与疲劳单模态模型的训练和调优。
    \item 多模态训练模块：完成 body+face 双输入、多任务输出模型的设计与训练。
    \item 推理与评估模块：实现统一测试集上的公平对比和视频级推理分析。
    \item 界面展示模块：实现本地可视化推理界面，用于结果演示。
\end{enumerate}

\section{数据来源说明}
本文主要使用的数据来源包括：

\begin{itemize}
    \item 分心源域数据：State Farm 驾驶行为分类数据集；
    \item 分心目标域数据：DMD 数据集中 S1 会话；
    \item 疲劳源域数据：NTHU 疲劳检测数据集；
    \item 疲劳目标域数据：DMD 数据集中 S5 会话；
    \item 多模态数据：基于 DMD 中配对的 body 与 face 视频构建的多模态样本。
\end{itemize}

\chapter{单模态驾驶分心识别方法设计与实验}

\section{任务定义与数据组织}
在分心识别部分，本文首先采用 State Farm 数据集作为源域训练数据。该数据集包含多种典型驾驶行为类别，图像质量较高、样本量充足，适合作为驾驶行为分类模型的初始训练基础。为了进一步评估模型在真实车内环境中的可用性，本文引入 DMD 数据集中的 S1 会话作为目标域数据，并将其中与本文研究目标高度相关的类别映射为四类典型驾驶状态，即 safe\_drive、radio、drinking 和 talking\_to\_passenger。

在数据组织过程中，本文重点保留了语义边界较清晰、样本质量较高的目标域类别，而没有将全部原始动作类别都直接纳入训练。一方面，这有助于降低目标域标签映射时的歧义；另一方面，也使得后续误判分析更具针对性。最终，本文构建了源域训练数据、目标域验证数据以及目标域 holdout 数据，用于完成分心识别模型的训练、调优和泛化验证。

\section{基线模型与训练设置}
本文在分心单模态识别中采用 YOLO11-cls 作为基础分类模型。选择该模型的原因在于：其一，该模型结构较为轻量，便于在现有实验环境下快速迭代；其二，模型在图像分类任务中具有较好的特征提取能力，适合作为后续多模态 body 分支的初始化基础。

前期实验中，本文尝试了输入分辨率调整、增强策略变化以及训练轮次变化等设置，并逐步观察不同策略对目标域效果的影响。在实验过程中，本文发现驾驶分心分类与常规通用图像分类存在明显差异。具体而言，驾驶行为类别中往往包含明确的左右语义，例如不同方向的手部动作、不同位置的操作行为等。如果直接使用常见的水平翻转增强，将会破坏类别语义，从而给模型带来错误监督。基于这一观察，本文在后续训练中关闭了水平翻转和垂直翻转，而保留了适度的颜色扰动、尺度缩放和随机擦除等增强方式。

\section{目标域适配策略}
仅使用 State Farm 数据集训练得到的模型虽然在源域上可以取得较高精度，但在 DMD S1 真实车内场景上表现并不稳定。其根本原因在于，源域与目标域在摄像视角、驾驶员姿态、光照条件和车辆环境等方面存在明显域差。因此，本文进一步构建了混合微调数据集，将经过清洗和映射的 DMD 样本与源域样本联合用于训练，以实现目标域适配。

在适配策略上，本文没有采用大步长重新训练，而是采用了较保守的微调方案。具体做法包括：以源域训练较优权重作为初始化，减小学习率，缩短训练轮次，并使用较弱的数据增强。实验表明，这种保守微调方式能够在不显著破坏源域特征表达能力的前提下，有效将模型决策边界向真实车内场景分布靠近。经过多轮实验比较后，本文最终确定 train24 所对应的权重作为当前分心单模态最佳模型。

\section{实验现象与结果分析}
从已有实验结果来看，分心单模态训练过程呈现出较为明确的规律。首先，在不考虑目标域适配时，模型在公开数据集上的性能较高，但在 DMD S1 上会出现明显的类别混淆，尤其是 safe\_drive 与其他相近行为之间的边界不够稳定。其次，在增强策略方面，简单增加增强强度并不能稳定提升目标域效果，反而有可能因为破坏姿态语义而造成性能下降。最后，在目标域混合微调后，模型在 DMD holdout 上的表现得到明显改善，说明目标域小样本校准在该任务中是必要且有效的。

进一步分析可知，body 模态对于分心识别的贡献主要体现在整体姿态、手部位置和中控交互等方面。相比之下，face 模态在该任务中提供的增益相对有限。这一结论为后续多模态模型设计提供了重要依据，即在分心任务中 body 分支应作为主导模态，而 face 分支更适合作为补充信息来源。

\section{阶段性结论}
综上，分心单模态部分可以得到以下几点结论：其一，驾驶行为分类任务中的增强策略需要考虑类别语义，尤其应避免使用会破坏左右动作含义的翻转增强；其二，源域模型难以直接适应真实车内场景，目标域混合微调是提升实际效果的关键步骤；其三，body 模态已经能够较好地表征驾驶分心状态，是后续多模态模型中分心分支的主要信息来源。

\chapter{单模态驾驶疲劳识别方法设计与实验}

\section{疲劳任务定义与数据集构建}
在疲劳识别部分，本文主要采用 face 模态作为单模态输入。其原因在于，疲劳状态通常通过眼睛开合、嘴部张开、打哈欠以及面部困倦表情等细粒度特征体现，而这些信息主要集中在面部区域。基于此，本文首先从 DMD S5 会话中整理出一个面向真实车内场景的小规模 face 疲劳数据集，并将样本划分为 fatigue 与 non\_fatigue 两类。

在标签构建时，本文没有将所有细粒度面部状态直接纳入训练，而是选择更加稳定的疲劳定义。例如，将明显的打哈欠和持续时间较长的闭眼片段归为 fatigue，将正常睁眼状态归为 non\_fatigue，而未将 opening、closing、blinking 和 undefined 等过渡状态直接纳入第一版训练。这种处理的主要考虑是：第一阶段目标是建立一个标签定义清晰、可稳定收敛的疲劳识别 baseline，而不是一开始就引入大量边界模糊状态。

\section{源域训练与跨域问题}
由于 DMD S5 数据规模有限，仅依赖小规模目标域训练难以获得稳定模型，因此本文进一步引入了 NTHU 疲劳数据集构建大规模源域训练集。基于该数据集训练得到的 face 疲劳分类模型在源域内部取得了较高准确率，说明模型本身能够有效学习典型疲劳表征。

然而，当该模型直接迁移到 DMD S5 目标域上时，本文观察到了显著的域偏移问题。最典型的现象是：模型在目标域测试集上几乎将所有样本都预测为 non\_fatigue，说明源域学到的特征表示并不能直接适应真实车内场景中的光照、视角和疲劳表现形式。与分心识别相比，疲劳识别对数据采集风格更敏感，这也是本文在疲劳线中必须重点处理的核心问题。

\section{目标域适配与混合微调}
为提升 face 疲劳模型在真实车内场景中的适用性，本文构建了 NTHU 与 DMD S5 的混合微调数据集，并以源域较优权重为初始化进行目标域适配。相较于直接使用大学习率重新训练，本文采用了较保守的微调方案，即降低学习率、减弱增强、缩短训练轮次，从而尽量保留源域模型的基础表征能力，同时让模型逐步适应目标域样本分布。

实验结果表明，经过混合适配后，模型在 S5 目标域上的性能得到了明显改善。进一步比较不同适配策略可以发现，一般性的目标域混合微调是有效的，但继续在 face 单模态上通过定向重加权方式强化某些疲劳子类型，并未带来稳定收益。这说明 face 单模态识别已经逐渐逼近其当前结构上限。最终，本文确定 train29 所对应的权重作为当前疲劳单模态最佳模型。

\section{误判分析与任务瓶颈}
通过对 S5 目标域误判样本的分析，本文发现 face 单模态模型的主要困难并不在于识别显著疲劳状态，而在于识别带遮挡或表征不充分的疲劳样本。例如，在打哈欠过程中若嘴部被手部部分遮挡，或者驾驶员处于轻度疲劳、眼部变化不明显的状态时，模型容易将 fatigue 误判为 non\_fatigue。相反，对于明显的闭眼或无遮挡打哈欠样本，模型通常能够给出较高置信度的正确预测。

这一现象说明，face 模态虽然非常适合建模疲劳状态，但在面对遮挡、角度变化和弱表征样本时仍然存在天然局限。这种瓶颈并不能单纯依靠继续在 face 单模态上堆叠训练技巧来完全解决，而更适合通过引入额外模态进行补充。也正因如此，本文后续引入 body 模态，希望利用驾驶姿态、手部位置等信息对 face 模态进行补充。

\section{阶段性结论}
综上，疲劳单模态部分的主要结论如下：其一，疲劳识别任务比驾驶分心识别更容易受到源域与目标域差异的影响；其二，单纯依赖大规模公开疲劳数据集训练不足以保证真实车内场景下的性能，目标域适配是必要步骤；其三，face 单模态模型对明显疲劳状态识别效果较好，但对遮挡型和弱表征疲劳样本仍存在明显瓶颈，这为后续引入 body 模态进行多模态融合提供了直接动机。

\chapter{双模态多任务模型设计与实验}

\section{设计动机}
在完成分心单模态和疲劳单模态建模之后，本文进一步考虑 body 与 face 两种模态之间的互补关系。body 模态更适合表达驾驶姿态、手部动作和整体行为，face 模态更适合表达眼睛状态、嘴部张开和疲劳表情。基于此，本文设计了 body+face 双模态多任务模型，同时输出分心识别结果和疲劳识别结果。

\section{数据组织方式}
本文构建了多模态数据集 dataset\_multimodal\_driver\_v1，其中每一条样本包含一张 body 图像、一张 face 图像以及对应的分心标签或疲劳标签。由于现实中并非每条样本同时具备完整的两类标签，因此数据集中使用 $-1$ 表示缺失标签，在训练时通过 masked loss 忽略对应损失项。该数据集中，S1 样本主要提供分心监督，S5 样本主要提供疲劳监督。

\section{网络结构}
多模态模型整体由 body encoder、face encoder、特征投影层、融合层以及两个任务头构成。body encoder 使用分心单模态模型训练得到的权重进行初始化，face encoder 使用疲劳单模态模型训练得到的权重进行初始化。两种模态的特征在融合层中进行拼接与映射，随后分别送入 distraction head 与 fatigue head 输出两个任务的预测结果。

\section{训练策略}
为了避免在多模态训练初期破坏单模态已经学到的有效特征，本文采用了两阶段训练策略：首先冻结两个编码器，仅训练融合层与任务头；随后在较小学习率下对整个模型进行联合微调。此外，在第二版多模态模型中，本文进一步提高了 fatigue 任务损失权重，以增强模型对疲劳任务的关注。

\section{实验结果与分析}
目前的对比实验表明，多模态模型在两个任务上均取得了较好效果。与单模态模型相比，双模态多任务模型在分心识别与疲劳识别上均有提升，尤其在疲劳识别任务上提升更为明显。多模态第二版模型在同一测试集上的结果如表~\ref{tab:current-results} 所示。

\begin{table}[H]
    \centering
    \caption{当前阶段主要模型结果对比}
    \label{tab:current-results}
    \begin{tabular}{>{\raggedright}p{4.2cm}cc}
        \toprule
        模型 & 分心准确率 & 疲劳准确率 \\
        \midrule
        分心单模态最佳模型 & 0.9474 & -- \\
        疲劳单模态最佳模型 & -- & 0.6795 \\
        多模态模型 v1 & 0.9474 & 0.8526 \\
        多模态模型 v2 & 0.9825 & 0.8590 \\
        \bottomrule
    \end{tabular}
\end{table}

从现有结果看，body+face 双模态模型在保持分心识别能力的同时，显著提升了 fatigue 任务表现，说明该模型能够利用模态间互补信息弥补单模态缺陷。

\chapter{系统实现与界面设计}

\section{视频对齐与推理流程}
为了保证多模态模型推理时两个输入流在时间上保持一致，本文实现了双视频对齐脚本，根据 DMD 标注文件中的 frame\_shift 信息对 body 和 face 视频进行同步修正。随后实现了统一的视频推理脚本，能够读取对齐后的双视频输入，调用多模态模型完成逐帧推理，并输出可视化视频与 csv 结果。

\section{本地化界面}
在模型训练和推理脚本基础上，本文进一步设计并实现了基于 tkinter 的本地界面原型。该界面目前已经支持以下功能：

\begin{enumerate}
    \item 选择多模态模型配置文件与权重文件；
    \item 选择 body 视频与 face 视频；
    \item 设置推理参数，如采样间隔、平滑窗口大小和同步偏移；
    \item 实时显示推理画面与当前分心、疲劳预测结果；
    \item 保存可视化视频与逐帧预测 csv。
\end{enumerate}

当前界面已经具备基础演示能力，为后续答辩展示与系统集成提供了支撑。

\chapter{总结与展望}

\section{全文总结}
本文围绕真实车内场景下的驾驶员分心与疲劳识别问题，完成了从单模态建模到多模态融合，再到本地系统实现的完整研究流程。首先，本文完成了分心单模态模型的源域训练与目标域适配；其次，完成了基于 face 模态的疲劳识别模型构建及目标域适配；随后，设计了 body+face 双模态多任务模型，并通过统一测试集上的公平对比验证了多模态模型相较于单模态模型的有效性；最后，实现了本地化推理界面原型，使研究结果具备了初步工程展示价值。

\section{后续工作展望}
尽管本文目前已经取得了阶段性成果，但仍有进一步优化空间，主要包括以下几个方面：

\begin{enumerate}
    \item 扩充更多真实车内场景的 body+face 配对样本，以进一步提升多模态模型的泛化能力；
    \item 引入更具针对性的融合结构与结构消融实验，以增强论文方法创新性；
    \item 引入时序建模方法，进一步减少视频推理中的短时抖动和误报现象；
    \item 继续完善本地界面交互与结果展示形式，形成更完整的论文演示系统。
\end{enumerate}

\end{document}
